\documentclass[10pt, letterpaper]{article}
\usepackage[ignoreheadfoot,top=2cm,bottom=2cm,left=2cm,right=2cm,footskip=0.8cm]{geometry}
\usepackage{fontawesome5}
\usepackage{hyperref}
\hypersetup{colorlinks=true,urlcolor=black}
\usepackage[french]{babel}
\usepackage[T1]{fontenc}
\usepackage{lmodern}
\pagestyle{empty}
\setlength{\parindent}{0pt}
\setlength{\parskip}{0.5em}
\begin{document}
\begin{center}
\textbf{\LARGE Objet : Candidature -- Quantitative Developer (Commodities \& Risk)}\\[0.2cm]
\textbf{\large Charly-Romy TANGA}\\[0.2cm]
\href{mailto:charlyromytanga.cytech@gmail.com}{\faEnvelope\ charlyromytanga.cytech@gmail.com} \quad$\cdot$\quad
\faPhone\ +33 6 15 42 25 74 \quad$\cdot$\quad
\href{https://www.linkedin.com/in/charly-romy-tanga}{\faLinkedin\ charly-romy.tanga} \quad$\cdot$\quad
\href{https://github.com/charlyromytanga}{\faGithub\ charlyromytanga}
\end{center}
Madame, Monsieur,

Je souhaite orienter ma carrière vers la modélisation quantitative appliquée aux marchés de l'énergie et des matières premières.

Mon expérience dans le secteur énergétique m'a permis de travailler sur des données nationales à grande échelle et de comprendre l'importance d'une analyse quantitative robuste pour les décisions stratégiques.

En parallèle, je développe des projets de modélisation du risque (VaR, Monte-Carlo) et de pricing d'options en C++/Python afin de renforcer mes compétences en finance quantitative appliquée aux marchés volatils.

Je souhaite rejoindre un environnement de trading ou d'analyse quantitative où je pourrai contribuer à la modélisation des risques et des dynamiques de marché des commodities.

Cordialement,\\
Charly-Romy TANGA
\end{document}
